\section{光的偏振}
\sourceproblem{选择1}{50}{56}
\problemtext 自然单色光照双缝，在两缝后各放一片偏振片，判断条纹变化。
\workings 若两偏振片透振方向未给定，不能唯一判定。平行时两路强度各减半、间距不变；正交时无干涉调制；一般相干项乘以两偏振方向内积，且还依入射偏振相关性。
\crosscheck 原图无透振方向条件，不能无条件断言条纹消失或只变暗。
\sourceproblem{选择2}{50}{56}
\problemtext 自然光和线偏振光非相干混合，通过转动检偏器后最大与最小强度比7，求原自然光与线偏振光强比。
\workings $I_{\max}=I_n/2+I_p$，$I_{\min}=I_n/2$，故 $I_p=3I_n$，$I_n/I_p=1/3$，选B。
\crosscheck 最小值来自自然光分量；采用非相干强度相加。
\sourceproblem{选择3}{50}{56}
\problemtext 强度 $I_0$ 的线偏振光经两片偏振片，透振方向与原偏振方向分别成 $\alpha$ 和90度，求出射光强。
\workings $I=I_0\cos^2\alpha\cos^2(90^\circ-\alpha)=I_0\sin^2(2\alpha)/4$，选A。
\crosscheck $\alpha=45^\circ$ 时达到 $I_0/4$；不能对初始线偏振光额外乘1/2。
\sourceproblem{选择4}{50}{56}
\problemtext 自然光强 $I_0$，两偏振片透振方向相夹30度，求出射强度。
\workings $I=(I_0/2)\cos^230^\circ=3I_0/8$，选D。
\crosscheck 第一片将理想自然光强减半，第二片用马吕斯定律。
\sourceproblem{选择5}{50}{56}
\problemtext 两片偏振片最初正交，自然光正入射，无光通过；将其中一片转过90度，判断光强。
\workings 转角 $\beta$ 从0增至90度时，$I=(I_0/2)\sin^2\beta$ 单调增大，选A。
\crosscheck 最终透振方向平行，透射为 $I_0/2$。
\sourceproblem{选择6}{50}{56}
\problemtext 自然光以60度由空气入射透明介质，反射光为线偏振光，判断折射光。
\workings 布儒斯特角 $i_B=60^\circ$，反射与折射方向正交，折射角30度。通常折射光仍含两种偏振分量，是部分偏振光，选C。
\crosscheck 折射率比为 $\tan60^\circ$；不能因反射全偏振就推断折射也全偏振。
\sourceproblem{选择7}{51}{57}
\problemtext 自然光以布儒斯特角由空气入射透明介质，判断反射光偏振方向。
\workings 理想非吸收介质界面在布儒斯特角的p偏振反射系数为零，故反射光为电矢量垂直入射面的线偏振光，选C。
\crosscheck 与折射光的偏振状态不同。
\sourceproblem{填空1}{51}{57}
\problemtext 转动检偏器区分自然光、线偏振光与部分偏振光。
\workings 理想自然光透射强度不变；线偏振光强度周期变化并有全暗；自然光与线偏振光的非相干混合光周期变化但不全暗。
\crosscheck 仅用转动线检偏器不能全面区分任意椭圆偏振与部分偏振，圆偏振也可给恒定强度。
\sourceproblem{填空2}{51}{57}
\problemtext 解释马吕斯定律 $I=I_0\cos^2\alpha$ 中各量。
\workings $I$ 为透过检偏器的线偏振光强，$I_0$ 为入射线偏振光强，$\alpha$ 为入射电矢量方向与检偏器透振方向的夹角。
\crosscheck 自然光入射时先求其经起偏器后的强度。
\sourceproblem{填空3}{51}{57}
\problemtext 自然光通过两片理想偏振片，夹角从30度变为45度，求转前与转后光强之比。
\workings $I_{30}/I_{45}=\cos^230^\circ/\cos^245^\circ=3/2$。
\crosscheck 若题目采用转后比转前则为2/3，注明比值顺序。
\sourceproblem{填空4}{51}{57}
\problemtext 自然光强 $I_0$，通过两偏振片后为 $I_0/4$，求透振方向夹角。
\workings $(I_0/2)\cos^2\alpha=I_0/4$，取轴线最小夹角得 $\alpha=45^\circ$。
\crosscheck 透振轴无正反之分，135度与45度等效。
\sourceproblem{填空5}{51}{57}
\problemtext 自然光由空气射向玻璃，折射角30度时反射光完全线偏振，求玻璃折射率。
\workings $i_B=90^\circ-30^\circ=60^\circ$，$n=\sin60^\circ/\sin30^\circ=\sqrt3$。
\crosscheck 与布儒斯特公式 $n=\tan i_B$ 一致。
\sourceproblem{填空6}{51}{57}
\problemtext 光由折射率 $n_1$ 的介质以布儒斯特角 $i_B$ 入射折射率 $n_2$ 的介质，求 $\tan i_B$。
\workings 对无吸收、通常非磁性介质，由 $r=90^\circ-i_B$ 和斯涅尔定律得 $\tan i_B=n_2/n_1$。
\crosscheck 折射率顺序不可颠倒；相同介质无反射，布儒斯特角退化。
